\chapter{1977年广西卷（理）}

\section{解答题}

\begin{problem}
\( x \) 为何值时：

\begin{enumerate}
\item \( \frac{|x|}{x} \) 的值是 \( 1 \)？
\item \( \frac{|x|}{x} \) 的值是 \( -1 \)？
\item \( \frac{|x|}{x} \) 没有意义？
\end{enumerate}
\end{problem}

\begin{solution}
根据分母 \(x\) 是否为零，分三种情形讨论：

\begin{enumerate}
\item 当 \(x>0\) 时，\(|x|=x\)，所以 \(\frac{|x|}{x}=1\)，故第（1）问的解为 \(x>0\).
\item 当 \(x<0\) 时，\(|x|=-x\)，所以 \(\frac{|x|}{x}=-1\)，故第（2）问的解为 \(x<0\).
\item 当 \(x=0\) 时分母为零，\(\frac{|x|}{x}\) 没有意义，故第（3）问的解为 \(x=0\).
\end{enumerate}
\end{solution}

\begin{problem}
不查表求：\(\lg2-2\lg4+\lg5+\lg16\) 的值.
\end{problem}

\begin{solution}
由对数的幂运算法则，\(\lg4=\lg\left(2^2\right)=2\lg2\)，\(\lg16=\lg\left(2^4\right)=4\lg2\). 又因为 \(\lg5=\lg\frac{10}{2}=\lg10-\lg2=1-\lg2\)，所以

\[
\lg2-2\lg4+\lg5+\lg16=\lg2-4\lg2+(1-\lg2)+4\lg2=1
\]
\end{solution}

\begin{problem}
计算：\(\cot\left( -60^\circ \right)\sin240^\circ-2\cos\left( -60^\circ \right)+\tan^2 150^\circ\).
\end{problem}

\begin{solution}
利用诱导公式，\(\cot\left( -60^\circ \right)=-\cot60^\circ=-\frac{\sqrt{3}}{3}\)，\(\sin240^\circ=-\sin60^\circ=-\frac{\sqrt{3}}{2}\)，\(\cos\left( -60^\circ \right)=\cos60^\circ=\frac{1}{2}\)，并且 \(\tan150^\circ=-\tan30^\circ=-\frac{\sqrt{3}}{3}\). 因此

\[
\begin{gathered}
\text{原式}=\left(-\frac{\sqrt{3}}{3}\right)\left(-\frac{\sqrt{3}}{2}\right)-2\cdot\frac{1}{2}+\left(-\frac{\sqrt{3}}{3}\right)^2\\
=\frac{1}{2}-1+\frac{1}{3}=-\frac{1}{6}
\end{gathered}
\]
\end{solution}

\begin{problem}
计算：\( \left( -1\frac{1}{2} \right)^2+81^{\frac{3}{4}}-\left( -1 \right)^0\div0.2+\left( -2 \right)^{-3} \).
\end{problem}

\begin{solution}
先逐项化简：\(\left(-1\frac{1}{2}\right)^2=\left(-\frac{3}{2}\right)^2=\frac{9}{4}\)，\(81^{\frac{3}{4}}=\left(3^4\right)^{\frac{3}{4}}=3^3=27\)，\(\left(-1\right)^0\div0.2=1\div\frac{1}{5}=5\)，以及 \(\left(-2\right)^{-3}=\frac{1}{\left(-2\right)^3}=-\frac{1}{8}\).

所以原式 \(=\frac{9}{4}+27-5-\frac{1}{8}=22+\frac{18}{8}-\frac{1}{8}=\frac{193}{8}=24\frac{1}{8}\).
\end{solution}

\begin{problem}
求 \(A\left(7,5\right)\) 和 \(B\left(1,-3\right)\) 两点间的距离.
\end{problem}

\begin{solution}
由两点间的距离公式，

\[
|AB|=\sqrt{\left(7-1\right)^2+\left(5-\left(-3\right)\right)^2}=\sqrt{6^2+8^2}=\sqrt{100}=10
\]

所以两点 \(A\left(7,5\right)\) 和 \(B\left(1,-3\right)\) 间的距离为 \(10\).
\end{solution}

\begin{problem}
红旗机械厂生产 \(360\text{ 台}\) 车床，原计划若干天完成，在党的十一大精神鼓舞下，实际每天比原计划多生产 \(12\text{ 台}\)，结果提前 \(8\text{ 天}\) 完成任务. 问实际每天生产多少台车床？
\end{problem}

\begin{solution}
设实际每天生产 \(x\text{ 台}\) 车床，则原计划每天生产 \(\left(x-12\right)\text{ 台}\). 由于原计划的日产量为正且实际日产量更多，必须有 \(x>12\). 按原计划需要 \(\frac{360}{x-12}\) 天，按实际计划需要 \(\frac{360}{x}\) 天.

由实际比原计划提前 \(8\) 天，得

\[
\frac{360}{x-12}-\frac{360}{x}=8
\]

因为 \(x>12\)，分母 \(x\left(x-12\right)\ne0\)，可以去分母，得到 \(360x-360\left(x-12\right)=8x\left(x-12\right)\)，即 \(8x^2-96x-4320=0\). 两边同除以 \(8\)，得

\(x^2-12x-540=0\)，\(\left(x-30\right)\left(x+18\right)=0\)

所以 \(x=30\) 或 \(x=-18\). 因为 \(x>12\)，舍去 \(x=-18\)，保留 \(x=30\). 代回检验时，原计划日产量为 \(30-12=18\) 台，原计划用时 \(360\div18=20\) 天，实际用时 \(360\div30=12\) 天，确实提前 \(20-12=8\) 天.

答：实际每天生产 \(30\text{ 台}\) 车床.
\end{solution}

\begin{problem}
在四边形 \(ABCD\) 中，\(E\)，\(F\) 分别是对角线 \(AC\)，\(BD\) 的中点，\(M\)，\(N\) 分别是 \(AB\)，\(CD\) 的中点，求证 \(EF\) 与 \(MN\) 互相平分.

\bitmapfigure[max width=5.6cm]{img/1977/guangxi_science/q7_fig1.png}
\end{problem}

\begin{solution}
连接 \(EM\)，\(EN\)，\(FM\)，\(FN\). 在三角形 \(ACD\) 中，\(E\) 和 \(N\) 分别是边 \(AC\) 和 \(CD\) 的中点. 由三角形中位线定理，\(EN\parallel AD\)，且 \(EN=\frac{1}{2}AD\).

在三角形 \(ABD\) 中，\(M\) 和 \(F\) 分别是边 \(AB\) 和 \(BD\) 的中点. 由三角形中位线定理，\(MF\parallel AD\)，且 \(MF=\frac{1}{2}AD\). 因此 \(MF\parallel EN\) 且 \(MF=EN\). 四边形 \(MFNE\) 有一组对边平行且相等，所以是平行四边形.

在平行四边形 \(MFNE\) 中，\(MN\) 和 \(FE\) 正好是两条对角线. 平行四边形的两条对角线互相平分，所以 \(EF\) 与 \(MN\) 互相平分.
\end{solution}

\begin{problem}
已知正四棱锥底面边长为 \(6\text{ 厘米}\)，侧棱是高的 \(2\) 倍. 求这棱锥的体积. （精确到 \(0.1\text{ 立方厘米}\)）.

\bitmapfigure[max width=4.8cm]{img/1977/guangxi_science/q8_fig1.png}
\end{problem}

\begin{solution}
设正四棱锥为 \(V\text{-}ABCD\)，底面中心为 \(O\)，高为 \(VO\). 因为底面 \(ABCD\) 是边长为 \(6\text{ 厘米}\) 的正方形，且 \(O\) 是其对角线交点，正方形的两条对角线互相垂直且平分，所以在直角三角形 \(AOD\) 中，\(AO=DO\)，\(\angle AOD=90^\circ\)，从而

\(2AO^2=AD^2=6^2\)，\(AO^2=18\)

在直角三角形 \(VOA\) 中，侧棱是高的 \(2\) 倍，即 \(VA=2VO\)，由勾股定理有 \(VA^2=VO^2+AO^2\). 因此 \(4VO^2=VO^2+18\)，解得 \(VO^2=6\)，由于高为正，\(VO=\sqrt{6}\text{ 厘米}\).

棱锥的底面积为 \(6^2=36\text{ 平方厘米}\)，所以体积为

\[
V_{\text{棱锥}}=\frac{1}{3}\times36\times\sqrt{6}=12\sqrt{6}\text{ 立方厘米}\approx29.3939\text{ 立方厘米}
\]

精确到 \(0.1\text{ 立方厘米}\) 时，\(29.3939\ldots\) 应四舍五入为 \(29.4\).

答：这棱锥的体积约为 \(29.4\text{ 立方厘米}\).
\end{solution}

\begin{problem}
把一块钢板冲成上面是半圆形，下面是矩形的零件，已知零件的周长是 \(P\)，怎样才能使冲出的零件面积最大？

\bitmapfigure[max width=6.0cm]{img/1977/guangxi_science/q9_fig1.png}
\end{problem}

\begin{solution}
设半圆的直径（也是矩形的宽）为 \(x\)，矩形的高为 \(h\)，零件面积为 \(S\). 半圆的弧长为 \(\frac{\pi x}{2}\)，所以周长条件给出 \(\frac{\pi x}{2}+x+2h=P\)，即

\[
h=\frac{1}{2}\left(P-\frac{\pi x}{2}-x\right)=\frac{2P-\left(\pi+2\right)x}{4}
\]

由于周长 \(P>0\)，且 \(x>0\) 与 \(h>0\)，变量范围为 \(0<x<\frac{2P}{\pi+2}\). 半圆面积为 \(\frac{\pi}{8}x^2\)，矩形面积为 \(xh\)，因此

\[
\begin{gathered}
S=\frac{\pi}{8}x^2+x\cdot\frac{2P-\left(\pi+2\right)x}{4}\\
=-\frac{\pi+4}{8}x^2+\frac{P}{2}x\\
=-\frac{\pi+4}{8}\left(x-\frac{2P}{\pi+4}\right)^2+\frac{P^2}{2\left(\pi+4\right)}
\end{gathered}
\]

因为 \(\pi+4>0\)，平方项前的系数为负，所以当 \(x=\frac{2P}{\pi+4}\) 时面积取得最大值. 这个值满足可行范围，因为 \(\pi+4>\pi+2\). 此时矩形的高为

\[
h=\frac{2P-\left(\pi+2\right)\frac{2P}{\pi+4}}{4}=\frac{P}{\pi+4}
\]

于是矩形的宽与高之比为 \(x:h=\frac{2P}{\pi+4}:\frac{P}{\pi+4}=2:1\).

答：当矩形的宽与高之比为 \(2:1\) 时，冲出的零件面积最大.
\end{solution}

\section{选作题}

\begin{problem}
解方程组 \(\begin{cases}
\sqrt[3]{\left(x+y\right)^2}+2\sqrt[3]{\left(x-y\right)^2}=3\sqrt[3]{x^2-y^2},\\
3x-2y=13
\end{cases}\).
\end{problem}

\begin{solution}
设 \(a=\sqrt[3]{x+y}\)，\(b=\sqrt[3]{x-y}\). 因为 \(x^2-y^2=\left(x+y\right)\left(x-y\right)\)，且实数的立方根满足乘法关系，所以
\(\sqrt[3]{\left(x+y\right)^2}=a^2\)，\(\sqrt[3]{\left(x-y\right)^2}=b^2\)，\(\sqrt[3]{x^2-y^2}=ab\). 原方程组等价于

\[
\begin{cases}
a^2+2b^2=3ab,\\
3x-2y=13
\end{cases}
\]

将第一式移项并因式分解，得 \(a^2-3ab+2b^2=0\)，即 \(\left(a-b\right)\left(a-2b\right)=0\). 所以需要分别求解下面两种情形.

\begin{enumerate}
\item 当 \(a-b=0\) 时，\(a=b\). 立方根函数在实数上单调递增，因此 \(x+y=x-y\)，从而 \(y=0\). 代入第二个方程得 \(3x=13\)，所以 \(x=\frac{13}{3}\). 这一情形给出 \(\left(x,y\right)=\left(\frac{13}{3},0\right)\).
\item 当 \(a-2b=0\) 时，\(a=2b\). 两边立方得 \(x+y=8\left(x-y\right)\)，即 \(7x-9y=0\). 联立第二个方程，得

\[
\begin{cases}
7x-9y=0,\\
3x-2y=13,
\end{cases}
\qquad y=\frac{7x}{9},
\qquad 3x-\frac{14x}{9}=13,
\qquad x=9,\qquad y=7
\]

得到 \(\left(x,y\right)=\left(9,7\right)\). 直接代入核对：当 \(\left(x,y\right)=\left(\frac{13}{3},0\right)\) 时 \(a=b\)；当 \(\left(x,y\right)=\left(9,7\right)\) 时 \(x+y=16\)，\(x-y=2\)，故 \(a=\sqrt[3]{16}=2\sqrt[3]{2}=2b\). 两组解都满足 \(3x-2y=13\)，所以没有遗漏.
\end{enumerate}

答：原方程组的解为 \(\left(x,y\right)=\left(\frac{13}{3},0\right)\) 或 \(\left(x,y\right)=\left(9,7\right)\).
\end{solution}

\begin{problem}
线段 \(AB\) 之三等分点依次为 \(C\)，\(D\)，以 \(CD\) 为直径画圆，在圆周上任取一点 \(P\). 求证 \(\tan\angle APC\cdot\tan\angle BPD=\frac{1}{4}\).

\bitmapfigure[max width=8.0cm]{img/1977/guangxi_science/q11_fig1.png}
\end{problem}

\begin{solution}
以下取 \(P\notin\{C,D\}\)，否则题中的正切无意义. 设 \(AC=CD=DB=a\)，于是 \(AD=2a\)，\(BC=2a\). 因为 \(CD\) 是所作圆的直径，且 \(P\) 在圆周上，由圆周角定理得 \(\angle CPD=90^\circ\)，即 \(PC\perp PD\).

过 \(D\) 作 \(DE\parallel CP\)，交 \(PB\) 于 \(E\). 在三角形 \(BCP\) 中，\(D\) 在 \(BC\) 上，\(E\) 在 \(BP\) 上，且 \(DE\parallel CP\)，所以三角形 \(BDE\) 与三角形 \(BCP\) 相似. 相似比为
\(\frac{BD}{BC}=\frac{a}{2a}=\frac{1}{2}\)，从而 \(DE=\frac{1}{2}CP\). 又因 \(DE\parallel CP\) 且 \(CP\perp PD\)，所以 \(DE\perp PD\)，三角形 \(PDE\) 是直角三角形，故
\(\tan\angle BPD=\frac{DE}{PD}\).

过 \(C\) 作 \(CF\parallel DP\)，交 \(PA\) 于 \(F\). 在三角形 \(ADP\) 中，\(C\) 在 \(AD\) 上，\(F\) 在 \(AP\) 上，且 \(CF\parallel DP\)，所以三角形 \(ACF\) 与三角形 \(ADP\) 相似. 相似比为
\(\frac{AC}{AD}=\frac{a}{2a}=\frac{1}{2}\)，从而 \(CF=\frac{1}{2}DP\). 又因 \(CF\parallel DP\) 且 \(DP\perp PC\)，所以 \(CF\perp PC\)，三角形 \(PCF\) 是直角三角形，故 \(\tan\angle APC=\frac{CF}{PC}\).

因此

\[
\tan\angle APC\cdot\tan\angle BPD=\frac{CF}{PC}\cdot\frac{DE}{PD}=\frac{\frac{1}{2}DP}{PC}\cdot\frac{\frac{1}{2}CP}{PD}=\frac{1}{4}
\]

证毕.
\end{solution}
